\section{Violence and Revolution}

\begin{quotex}
The cause of all wars and revolutions --- in a word, of all violence --- is always the same: the \emph{negation of hierarchy}.
\flright{\textsc{Valentin Tomberg}}

What is needed is not a revolution in the opposite direction, but the opposite of a revolution.

\flright{\textsc{Joseph de Maistre}}

While the world does not understand these three words: Truth, Reason, Justice, and these: Duty, Hierarchy, Society, the revolutionary motto, `Liberty, Equality, Fraternity,' will be nothing but a threefold lie.
\flright{\textsc{Eliphas Levi}}
\end{quotex}


Little comment is needed, as these three Hermetists speak clearly. If revolution is the negation of hierarchy, then the opposite of a revolution is the restoration of hierarchy. Eliphas Levi proposes the motto for the counter-revolution based on Truth, Reason and Justice. Eliphas describes the coming government:

\begin{quotex}
The government of the future will be that whose model is shown to us in nature, by the family, and in the religious world by the pastoral hierarchy.
\end{quotex}


This is recognized more so by initiates:

\begin{quotationx}
The idea of a perfect and immutable order in nature, the notion of an ascending hierarchy and of a descending influence in all beings, had furnished to the ancient hierophants the first classification of the whole of natural history.

The initiate respects the hierarchy before all, he loves and preserves order ... 
\end{quotationx}


Valentin Tomberg traces this to the founding myths of the West

\begin{quotex}
The story of Cain and Abel is a \emph{myth}, i.e., it expresses, under the form of a tale of a particular case, an ``eternal'' idea. Consequently, it pertains to \emph{time}, to \emph{history}, and not to \emph{space} and its structure. It shows us how brothers can become mortal enemies through the very fact that they worship the same God in the same way. The source of religious wars is revealed here: and it is not the difference in dogma nor forms of worship or ritual which is their cause, but exclusively the pretention to equality, or if one prefers, the \emph{negation of hierarchy}. This is also the world's first revolution --- the archetype of all revolutions which have taken place and which will take place in the future of humanity. For the cause of all wars and revolutions --- in a word, of all violence --- is always the same: the negation of hierarchy.
\end{quotex}



\flrightit{Posted on 2010-10-13 by Cologero}

\begin{center}* * *\end{center}

\begin{footnotesize}
\begin{sffamily}
\texttt{Count Cagliostro on 2010-10-13 at 06:07 said:}

Very insightful!

But still, where is the point where respect to hierarchy becomes a tool so that the elites preserve their -mostly materialistic- privileges?

\hfill

\texttt{Ernest on 2010-10-13 at 06:43 said:}

I'm a little confused about how hierarchy is supposed to be negated in the myth of Cain and Abel.

Is it simply that the younger brother is favoured over the elder?

Or is it that by favouring Abel, God elevates his position above his older brother, and Cain pretends to equality and revolts against God's judgment by killing his promoted brother?

In his On Being a Pagan, Alain de Benoist notes:

``As a farmer, Cain is by that status rooted, attached to the soil that Yahweh has cursed because of Adam... He has chosen, as Levinas puts it, Totality as opposed to Infinity, the ``pagan'' conquest of space against the Hebrew possession of time as eternity. For this attachment to a given soil, rootedness, bears within itself the signs of everything the Bible stigmatizes as idolatry. These include the distinct cities, patriotism, the state and reasons of state, the frontier that distinguishes citizen from foreigner, the vocation of the soldier, politics, and so forth. Whereas by his sacrifice, Abel shows he keeps his spirit totally open for Yahweh, Cain's sacrifice asks God to sanctify the kind of existence that has earned God's disapproval because it is a manifestation of the increased autonomy man is seeking. Just like Adam, Cain reveals his pride, and this is why he is condemned. In fact, the principle cause of Cain's condemnation is not Abel's murder, but Cain's refusal to humble himself by repenting.''

Cain of course, later founded the first city, Enoch and ``did not restrict himself to engendering urban civilization, the one where history is made, but he also forms the first link in a long chain of inventors of civilization. One of his descendants, Yubal, was the first musician. Another Tubal-Cain is the ancestor of smiths, and it is to him we owe the discovery of metallurgy. In this he is considered to be the first specialist in the art of war, a fact that of course earns him God's personal disapproval.''

\hfill

\texttt{Cologero on 2010-10-13 at 09:05 said:}

The system where an ``elite'' hold power for their own material benefit (``oligarchy''), can only be maintained through force and violence; it is thus unstable and will incite revolution. ``Synarchy'' is a system hierarchically organized based on spiritual depth, intelligence, or accomplishment that is intended for the benefit of all. (See, for example, \textit{City of the Sun}\footnote{\url{http://www.gornahoor.net/?p=518}}.)

People today are suspicious of the latter's claims to hierarchy based on spirituality or intelligence. Nevertheless, those same people readily accept differences based on wealth, athletic prowess, and physical beauty.

\hfill

\texttt{James O'Meara on 2010-10-13 at 14:40 said:}

I don't usually have much use for Ken Wilber, but in his Sex Ecology and Spirituality which I reviewed for David Fideler's old Alexandria he comes up with the term `holoachy' composed of holons, of course, who are inherently interconnected in which the higher levels are naturally ruling for the betterment of the lower, rather than `hierarchy' in which the former force themselves on the latter for their own selfish goals. He agrees with your comment about needing to address the issue of historical cases of exploitation as motives for denying or denouncing hierarchic views. I don't recall any examples but something along the lines of a natural father looking out for his children, verus a Fagin-like exploiter.

\hfill

\texttt{Will on 2010-10-14 at 08:38 said:}

It seems to me that the problem with the word ``hierarchy'' is the political baggage it has picked up over the years. But the original meaning of the word is ``sacred order,'' from the Greek hieros -- `sacred' or `holy' -- and arche -- which means not only `order' but also `origin' or `beginning.' So the word denotes the sacred, natural, and thus right order of the universe.

But this meaning has been lost, and now the common understanding of the word associates it with artificial and corrupt power structures. I like the idea conveyed by Wilber's alternative term, but in my view, it would be better to recover the original sense of ``hierarchy,'' (which is identical to what Wilber writes of) rather than to play the PC game of constantly modifying our vocabulary because current fashion imposes a new meaning onto a word.

\hfill

\texttt{zero is zero on 2010-10-14 at 09:31 said:}

Question: I once read a passage by Evola where he criticized anarchism as an ideology for pathologically weakening the higher principles within one's self (and not just higher principles in external society). I am paraphrasing of course. I can't find the passage though, although I remember it was on the internet. Cologero or anyone do you know the passage I am talking about?

\hfill

\texttt{Matt on 2010-10-14 at 17:58 said:}

Zero,

My guess is that the passage is either from Men Amongst the Ruins, or Pagan Imperialism (which Cologero has much of its contents posted on the site).

\hfill

\texttt{Will on 2010-10-14 at 20:19 said:}

zero,

I would check Ride the Tiger, and also the chapter on ``Youth, Beats, and Right Wing Anarchism'' from The Bow and the Club, which is available on the Evola As He Is website.
\url{http://thompkins\_cariou.tripod.com/id99.html}


My understanding of his opinion of anarchism is that, as a metaphysical position, anarchism (particularly the various political movements of that name) was anti-Traditional, because it was essentially anti-order and pro-chaos. However, in the modern, anti-Traditional world, in which the external laws are no longer connected to any legitimate authority, anarchism somewhat paradoxically becomes a viable position, provided that one has the proper inner orientation. One obeys an inner law, rather like Herodotus says of the Spartans.

I think his view also is related to Nietzsche's idea that a civilization teetering on the brink of collapse ought to be pushed, rather than pulled back.

\hfill

\texttt{James O'Meara on 2010-10-15 at 18:13 said:}

Will,

I agree; E's anti-bourgeois position welcomed threats to the established order from any direction, though he criticized their stated goals or lack thereof . Better to go `on the road' seeking `kicks' than pursue, say, an academic career.

Trevor Lynch has reviewed The Dark Knight as presenting traditionalist themes out of the mouths of `evil doers' such as the League of Shadows or the Joker, including ``a civilization teetering on the brink of collapse ought to be pushed League, Joker , rather than pulled back Batman, of course .

\hfill

\texttt{zero is zero on 2011-01-24 at 15:06 said:}

Thank you Mark and Will for replying. I have finally tracked down the passage I was looking for, and I was mistaken in attributing it to Evola. It is in fact by Nietzsche, quoted in Hannson's introduction to Men Among the Ruins:

``... There is every reason for comparing the Christian and the anarchist because the impulse of both is toward destruction... . ``The Christian and anarchist are both decadents, both are only able to act for the purpose of disintegrating, poi-soning, degrading, and blood-sucking. They both have the impulse of mortal ha-tred toward anything that stands up, and is great, and is lasting and shows promise for the future... '' (The Antichrist, aphorism 58)

Although I wonder if Nietzsche understood ``anarchist'' then, in the late 19th century, in the meaning that it acquired over the course of the last century.

\hfill

\texttt{zero is zero on 2011-01-24 at 15:21 said:}

In my previous post ``Mark'' should have read ``Matt''; Sorry for that.

\hfill

\texttt{Ron on 2012-02-21 at 16:47 said:}

Cologero, where can I find a Hermetic order in the spirit of Levi, Tomberg, et al, one that is both Catholic and esoteric?

\hfill

\texttt{Cologero on 2012-02-21 at 22:49 said:}

Ron, I doubt that any such group, if it exists, would advertise on the Internet and accept any and all applicants. In any case, as you may recall, Tomberg specifically requested that no organization or order be established based on his meditations ... not that this has deterred certain elements from doing so.

Nevertheless, he did expect some ``friends'' to gather loosely around his teachings, developing them further, particularly to provide meditations on the minor arcana. I am not aware of anyone doing that, certainly not publicly.

Inspired by Louis Claude de Saint-Martin, we created a ``Society of Independents'', or SI, which is also the abbreviation in French. SI could also stand for Superieurs Inconnus (Unknown Superiors), alluding to Unknown Friends. It has not yet found a spiritual centre to undertake such a project.


\hfill

\end{sffamily}
\end{footnotesize}

